\section{Descripción de la Solución}
La solución es un torniquete a escala que se abre solo cuando el estudiante
muestra un código válido desde su celular. El código es un cuadro de puntos
llamado código QR, como los que se leen con la cámara del teléfono.

\subsection{Cómo Funciona, Paso a Paso}
\begin{enumerate}
	\item El estudiante abre una página web desde su celular y ve su código QR.
	      El código cambia cada pocos segundos.
	\item Al llegar al torniquete, muestra la pantalla a una cámara pequeña
	      (el lector ESP32-CAM), que captura la imagen y lee el código.
	\item Una segunda placa electrónica (el ESP32) recibe el código y lo envía
	      por Wi-Fi al servidor, que es la computadora central del sistema.
	\item El servidor revisa si el código es auténtico y sigue vigente, y
	      responde si el estudiante puede pasar.
	\item Si la respuesta es afirmativa, el torniquete se abre y una luz verde
	      lo indica. Si no, una luz y un zumbador avisan que no puede pasar.
	\item El servidor guarda cada intento, aceptado o rechazado, con fecha y
	      hora en una base de datos, que es el archivo donde se almacena la
	      información.
	\item El personal de vigilancia y Control Escolar consultan esos registros
	      en un panel web: ingresos por hora, intentos rechazados y estudiantes
	      que llevan varios días sin entrar.
\end{enumerate}

\begin{figure}[tbp]
	\caption{Cómo Funciona el Torniquete}
	\label{fig:flujo}
	\centering
	\begin{tikzpicture}[
		node distance=0.9cm and 1.0cm,
		box/.style={draw, rounded corners, align=center, minimum height=1.2cm,
				minimum width=2.4cm, font=\footnotesize},
		arr/.style={-{Latex[length=2mm]}, thick}
		]
		\node[box] (cel) {Celular\\(muestra el QR)};
		\node[box, right=of cel] (cam) {ESP32-CAM\\(lee el código)};
		\node[box, right=of cam] (esp) {ESP32\\(consulta y abre)};
		\node[box, below=of esp] (srv) {Servidor y\\base de datos\\(validan y guardan)};
		\node[box, left=of srv] (panel) {Panel web\\(vigilancia y\\Control Escolar)};
		\draw[arr] (cel) -- (cam);
		\draw[arr] (cam) -- (esp);
		\draw[arr] (esp) -- (srv) node[midway, right, font=\scriptsize]{Wi-Fi};
		\draw[arr] (srv) -- (panel) node[midway, above, font=\scriptsize]{consulta};
	\end{tikzpicture}
	\figurenote{Elaboración propia. El servidor responde si el código es válido y el ESP32 libera el torniquete; cada intento queda registrado en la base de datos.}
\end{figure}

\subsection{Por Qué el Código Cambia}
Cada código lleva una firma digital llamada HMAC, una contraseña de un solo
uso que solo el servidor sabe calcular. Si alguien cambia el código, la firma
deja de coincidir y se rechaza. Como además el código se renueva cada pocos
segundos, una captura de pantalla compartida deja de servir al poco tiempo.
Este tipo de firma resiste el reenvío y la modificación del mensaje
\parencite{subpratatsavee2015hmac}, y la renovación periódica ya se probó en
una universidad del Reino Unido, donde un código renovado cada 20 segundos
impidió que los estudiantes se reenviaran capturas \parencite{nwabuwe2023dynamicqr}.

\subsection{Límite de la Solución}
El código que cambia impide reutilizar capturas, pero no impide que un
estudiante preste su celular a otra persona. Por eso el proyecto busca reducir
la suplantación, no eliminarla. Además, como la medición se hace en un solo
acceso y sin grupo de comparación, los resultados se tomarán como evidencia
orientativa, no como prueba definitiva.
